\begin{Exercise}[label={ex:simmetrica}]
    Si scriva una funzione che calcola la differenza simmetrica degli elementi di due liste ordinate 
    in senso crescente e prive di duplicati, restituendola come una nuova lista (allocata allo scopo), 
    anch’essa ordinata. La differenza simmetrica è costituita dagli elementi che appartengono a una delle due liste ma non all’altra
    lista (contiene cioè tutti gli elementi che non sono in comune alle due liste).
\end{Exercise}

\begin{Answer}[ref={ex:simmetrica}]
  	\lstinputlisting[style=cstylesmall]{src/simmetrica.c}
\end{Answer}